% ============================================================================
% DICCIONARIO DE DATOS — SADUTO
% Sistema de Analítica de Datos Utilizando Técnicas Open Source Intelligence
% Caso: Vicerrectorado de Grado — Escuela Militar de Ingeniería
% ============================================================================
\documentclass[11pt,a4paper]{article}

% ── Paquetes ──────────────────────────────────────────────────────────────
\usepackage[utf8]{inputenc}
\usepackage[spanish]{babel}
\usepackage[T1]{fontenc}
\usepackage{lmodern}
\usepackage[margin=2cm]{geometry}
\usepackage{graphicx}
\usepackage{xcolor}
\usepackage{titlesec}
\usepackage{fancyhdr}
\usepackage{hyperref}
\usepackage{enumitem}
\usepackage{booktabs}
\usepackage{longtable}
\usepackage{tabularx}
\usepackage{float}
\usepackage{tcolorbox}
\usepackage{array}
\usepackage{pdflscape}
\usepackage{multirow}
\usepackage{colortbl}

% ── Colores institucionales ───────────────────────────────────────────────
\definecolor{emiblue}{HTML}{0D47A1}
\definecolor{emiyellow}{HTML}{FFD700}
\definecolor{emidark}{HTML}{002171}
\definecolor{emigray}{HTML}{F5F5F5}
\definecolor{emilightblue}{HTML}{E3F2FD}
\definecolor{tableheader}{HTML}{0D47A1}
\definecolor{tablerow}{HTML}{F0F7FF}

% ── Configuración de hipervínculos ────────────────────────────────────────
\hypersetup{
  colorlinks=true,
  linkcolor=emiblue,
  urlcolor=emiblue,
  pdftitle={Diccionario de Datos — SADUTO},
  pdfauthor={Alvaro Santiago Encinas Flores},
}

% ── Formato de títulos ────────────────────────────────────────────────────
\titleformat{\section}{\LARGE\bfseries\color{emiblue}}{}{0em}{}[\titlerule]
\titleformat{\subsection}{\Large\bfseries\color{emidark}}{}{0em}{}
\titleformat{\subsubsection}{\large\bfseries\color{emiblue}}{}{0em}{}

% ── Encabezado y pie de página ────────────────────────────────────────────
\pagestyle{fancy}
\fancyhf{}
\fancyhead[L]{\small\color{emiblue}\textbf{SADUTO} — Diccionario de Datos}
\fancyhead[R]{\small\color{gray}v1.0}
\fancyfoot[C]{\thepage}
\fancyfoot[R]{\small\color{gray}EMI Bolivia — 2026}
\renewcommand{\headrulewidth}{1pt}
\renewcommand{\headrule}{\hbox to\headwidth{\color{emiblue}\leaders\hrule height \headrulewidth\hfill}}

% ── Macro para cabecera de tabla ──────────────────────────────────────────
\newcommand{\tableheadrow}{\rowcolor{tableheader}}
\newcommand{\tablehead}[1]{\textcolor{white}{\textbf{#1}}}
\newcommand{\pk}{\textcolor{emiblue}{\textbf{PK}}}
\newcommand{\fk}{\textcolor{orange!80!black}{\textbf{FK}}}
\newcommand{\nn}{\textcolor{red!70!black}{\textbf{NN}}}

% ── Caja informativa ─────────────────────────────────────────────────────
\newtcolorbox{infobox}[1][]{
  colback=emilightblue,
  colframe=emiblue,
  fonttitle=\bfseries,
  title=#1,
  rounded corners,
  boxrule=0.5pt
}

% ══════════════════════════════════════════════════════════════════════════
\begin{document}

% ── Portada ───────────────────────────────────────────────────────────────
\begin{titlepage}
\centering
\vspace*{2cm}

{\Huge\bfseries\color{emiblue}SADUTO\par}
\vspace{0.3cm}
{\large\color{emidark}Sistema de Analítica de Datos Utilizando\\Técnicas Open Source Intelligence\\Caso: Vicerrectorado de Grado\par}

\vspace{1cm}
\noindent\makebox[\textwidth]{\color{emiblue}\rule{0.6\textwidth}{2pt}}
\vspace{1cm}

{\LARGE\bfseries Diccionario de Datos\par}
\vspace{0.3cm}
{\large Catálogo técnico de la estructura de la base de datos\par}

\vfill

\begin{tabular}{ll}
\textbf{Autor:}         & Alvaro Santiago Encinas Flores \\
\textbf{Carrera:}       & Ing.\ de Sistemas — EMI 2026 \\
\textbf{Institución:}   & Escuela Militar de Ingeniería (EMI) \\
\textbf{Motor BD:}      & SQLite 3 \\
\textbf{Archivo:}       & \texttt{data/osint\_emi.db} \\
\textbf{Total tablas:}  & 30 tablas \\
\textbf{Versión:}       & 2.0 \\
\textbf{Fecha:}         & Mayo 2026 \\
\end{tabular}

\vspace{1cm}
\noindent\makebox[\textwidth]{\color{emiyellow}\rule{0.4\textwidth}{3pt}}
\end{titlepage}

% ── Tabla de contenidos ───────────────────────────────────────────────────
\tableofcontents
\newpage

% ══════════════════════════════════════════════════════════════════════════
\section{Introducción}

\subsection{Propósito}
Este documento cataloga la estructura completa de la base de datos del sistema \textbf{SADUTO} (Sistema de Analítica de Datos Utilizando Técnicas Open Source Intelligence — Caso: Vicerrectorado de Grado). Para cada tabla se especifica: nombre del campo, tipo de dato, restricciones de nulidad (\nn), claves primarias (\pk), claves foráneas (\fk) y una descripción funcional del campo.

\subsection{Convenciones}
\begin{table}[H]
\centering
\begin{tabularx}{\textwidth}{l l X}
\toprule
\textbf{Símbolo} & \textbf{Significado} & \textbf{Descripción} \\
\midrule
\pk & Primary Key & Identificador único de la fila. Auto-incremental. \\
\fk & Foreign Key & Referencia a la clave primaria de otra tabla. \\
\nn & Not Null & El campo no admite valores nulos. \\
\bottomrule
\end{tabularx}
\end{table}

\subsection{Diagrama de relaciones (resumen)}
Las principales relaciones entre tablas son:

\begin{itemize}
  \item \texttt{dato\_recolectado} $\rightarrow$ \texttt{fuente\_osint} (origen de cada dato).
  \item \texttt{dato\_procesado} $\rightarrow$ \texttt{dato\_recolectado} (dato limpio del crudo).
  \item \texttt{analisis\_sentimiento} $\rightarrow$ \texttt{dato\_procesado} (resultado de IA).
  \item \texttt{comentario} $\rightarrow$ \texttt{dato\_recolectado} (comentarios de un post).
  \item \texttt{comentario} $\rightarrow$ \texttt{fuente\_osint} (fuente del comentario).
  \item \texttt{analisis\_comentario} $\rightarrow$ \texttt{comentario} (sentimiento del comentario).
  \item \texttt{alerta} $\rightarrow$ \texttt{usuario} (asignación y resolución).
  \item \texttt{configuracion\_alerta} $\rightarrow$ \texttt{usuario} (quién creó la regla).
  \item \texttt{log\_actividad} $\rightarrow$ \texttt{usuario} (registro de acciones).
\end{itemize}

% ══════════════════════════════════════════════════════════════════════════
\section{Tablas del módulo OSINT — Recolección}

% ── Tabla: fuente_osint ───────────────────────────────────────────────────
\subsection{Tabla: fuente\_osint}
\textit{Configuración de las redes sociales y plataformas de origen de datos OSINT.}

\begin{longtable}{p{4.2cm} p{2.8cm} c c c p{4.5cm}}
\toprule
\tableheadrow
\tablehead{Campo} & \tablehead{Tipo} & \tablehead{PK} & \tablehead{FK} & \tablehead{NN} & \tablehead{Descripción} \\
\midrule
\endhead
id\_fuente & INTEGER & \pk & & & Identificador único de la fuente. \\
\rowcolor{tablerow}
nombre\_fuente & VARCHAR(100) & & & \nn & Nombre descriptivo (ej.\ ``EMI Oficial Facebook''). \\
tipo\_fuente & VARCHAR(50) & & & \nn & Tipo de plataforma: Facebook, TikTok, Web. \\
\rowcolor{tablerow}
url\_fuente & VARCHAR(255) & & & \nn & URL de la fuente. \\
identificador & VARCHAR(100) & & & & ID interno de la plataforma. \\
\rowcolor{tablerow}
puntuacion\_confiabilidad & DECIMAL(3,2) & & & & Confiabilidad de la fuente (0.00--1.00). Defecto: 0.80. \\
activa & BOOLEAN & & & & Estado activo/inactivo. Defecto: 1. \\
\rowcolor{tablerow}
fecha\_creacion & TIMESTAMP & & & & Fecha de registro de la fuente. \\
fecha\_ultima\_recoleccion & TIMESTAMP & & & & Última vez que se extrajeron datos. \\
\rowcolor{tablerow}
total\_registros\_recolectados & INTEGER & & & & Contador acumulado de datos extraídos. \\
\rowcolor{tablerow}
metadata\_json & TEXT & & & & Metadatos adicionales en formato JSON. \\
es\_oficial & BOOLEAN & & & & Si es fuente oficial de la EMI. Defecto: 0. \\
\bottomrule
\caption{Tabla \texttt{fuente\_osint} — Fuentes OSINT (4 registros)}
\end{longtable}

% ── Tabla: dato_recolectado ───────────────────────────────────────────────
\subsection{Tabla: dato\_recolectado}
\textit{Datos extraídos crudos de cada fuente, con fechas y métricas de interacción (engagement).}

\begin{longtable}{p{4.2cm} p{2.8cm} c c c p{4.5cm}}
\toprule
\tableheadrow
\tablehead{Campo} & \tablehead{Tipo} & \tablehead{PK} & \tablehead{FK} & \tablehead{NN} & \tablehead{Descripción} \\
\midrule
\endhead
id\_dato & INTEGER & \pk & & & Identificador único del dato. \\
\rowcolor{tablerow}
id\_fuente & INTEGER & & \fk & \nn & Referencia a \texttt{fuente\_osint}. \\
id\_externo & VARCHAR(100) & & & \nn & ID de la publicación en la plataforma original. \\
\rowcolor{tablerow}
fecha\_publicacion & TIMESTAMP & & & \nn & Fecha original de la publicación. \\
fecha\_recoleccion & TIMESTAMP & & & & Fecha de extracción por el scraper. \\
\rowcolor{tablerow}
contenido\_original & TEXT & & & \nn & Texto íntegro de la publicación. \\
autor & VARCHAR(100) & & & & Autor/cuenta que publicó. \\
\rowcolor{tablerow}
engagement\_likes & INTEGER & & & & Cantidad de ``me gusta''. Defecto: 0. \\
engagement\_comments & INTEGER & & & & Cantidad de comentarios. Defecto: 0. \\
\rowcolor{tablerow}
engagement\_shares & INTEGER & & & & Cantidad de veces compartido. Defecto: 0. \\
engagement\_views & INTEGER & & & & Cantidad de visualizaciones. Defecto: 0. \\
\rowcolor{tablerow}
tipo\_contenido & VARCHAR(30) & & & & Tipo: post, comentario, video, etc. \\
url\_publicacion & VARCHAR(500) & & & & URL directa a la publicación. \\
\rowcolor{tablerow}
metadata\_json & TEXT & & & & Datos adicionales en JSON. \\
procesado & BOOLEAN & & & & Indica si fue procesado por el ETL. Defecto: 0. \\
\rowcolor{tablerow}
fecha\_procesamiento & TIMESTAMP & & & & Fecha del procesamiento ETL. \\
\bottomrule
\caption{Tabla \texttt{dato\_recolectado} — Datos crudos (214 registros)}
\end{longtable}

% ── Tabla: comentario ─────────────────────────────────────────────────────
\subsection{Tabla: comentario}
\textit{Comentarios extraídos de publicaciones, con soporte para hilos (respuestas a otros comentarios).}

\begin{longtable}{p{4.2cm} p{2.8cm} c c c p{4.5cm}}
\toprule
\tableheadrow
\tablehead{Campo} & \tablehead{Tipo} & \tablehead{PK} & \tablehead{FK} & \tablehead{NN} & \tablehead{Descripción} \\
\midrule
\endhead
id\_comentario & INTEGER & \pk & & & Identificador único. \\
\rowcolor{tablerow}
id\_post & INTEGER & & \fk & \nn & Ref.\ a \texttt{dato\_recolectado} (post padre). \\
id\_fuente & INTEGER & & \fk & \nn & Ref.\ a \texttt{fuente\_osint}. \\
\rowcolor{tablerow}
id\_externo & VARCHAR(100) & & & & ID en la plataforma original. \\
autor & VARCHAR(100) & & & & Nombre del autor. \\
\rowcolor{tablerow}
contenido & TEXT & & & \nn & Texto del comentario. \\
fecha\_publicacion & TIMESTAMP & & & & Fecha de publicación. \\
\rowcolor{tablerow}
fecha\_recoleccion & TIMESTAMP & & & & Fecha de extracción. \\
likes & INTEGER & & & & Likes del comentario. Defecto: 0. \\
\rowcolor{tablerow}
respuestas & INTEGER & & & & Cantidad de respuestas. Defecto: 0. \\
es\_respuesta & BOOLEAN & & & & Si es respuesta a otro comentario. \\
\rowcolor{tablerow}
id\_comentario\_padre & INTEGER & & \fk & & Ref.\ a \texttt{comentario} (auto-referencia). \\
procesado & BOOLEAN & & & & Si fue analizado. Defecto: 0. \\
\rowcolor{tablerow}
metadata\_json & TEXT & & & & Metadatos adicionales. \\
\bottomrule
\caption{Tabla \texttt{comentario} — Comentarios (534 registros)}
\end{longtable}

% ══════════════════════════════════════════════════════════════════════════
\section{Tablas del módulo ETL — Procesamiento}

% ── Tabla: dato_procesado ─────────────────────────────────────────────────
\subsection{Tabla: dato\_procesado}
\textit{Datos limpios y enriquecidos tras pasar por el pipeline ETL. Incluye métricas temporales y de engagement normalizadas.}

\begin{longtable}{p{4.2cm} p{2.8cm} c c c p{4.5cm}}
\toprule
\tableheadrow
\tablehead{Campo} & \tablehead{Tipo} & \tablehead{PK} & \tablehead{FK} & \tablehead{NN} & \tablehead{Descripción} \\
\midrule
\endhead
id\_dato\_procesado & INTEGER & \pk & & & Identificador único. \\
\rowcolor{tablerow}
id\_dato\_original & INTEGER & & \fk & \nn & Ref.\ a \texttt{dato\_recolectado}. \\
contenido\_limpio & TEXT & & & \nn & Texto normalizado y limpio. \\
\rowcolor{tablerow}
longitud\_texto & INTEGER & & & & Longitud en caracteres. \\
cantidad\_palabras & INTEGER & & & & Número de palabras. \\
\rowcolor{tablerow}
fecha\_publicacion\_iso & TIMESTAMP & & & \nn & Fecha normalizada ISO 8601. \\
anio & INTEGER & & & & Año de publicación. \\
\rowcolor{tablerow}
mes & INTEGER & & & & Mes (1--12). \\
dia\_semana & INTEGER & & & & Día de la semana (0=lunes). \\
\rowcolor{tablerow}
hora & INTEGER & & & & Hora de publicación (0--23). \\
semestre & VARCHAR(10) & & & & Semestre académico (ej.\ ``2024-II''). \\
\rowcolor{tablerow}
es\_horario\_laboral & BOOLEAN & & & & Si fue publicado en horario laboral. \\
engagement\_total & INTEGER & & & & Suma de likes + comments + shares. \\
\rowcolor{tablerow}
engagement\_normalizado & DECIMAL(5,2) & & & & Engagement escalado 0--100. \\
ratio\_engagement & DECIMAL(8,4) & & & & Engagement / alcance. \\
\rowcolor{tablerow}
categoria\_preliminar & VARCHAR(50) & & & & Categoría inicial asignada. \\
idioma\_detectado & VARCHAR(10) & & & & Idioma. Defecto: ``es''. \\
\rowcolor{tablerow}
contiene\_mencion\_emi & BOOLEAN & & & & Si menciona a la EMI. Defecto: 0. \\
sentimiento\_basico & VARCHAR(20) & & & & Sentimiento por reglas simples. \\
\rowcolor{tablerow}
fecha\_procesamiento & TIMESTAMP & & & & Fecha del procesamiento ETL. \\
version\_etl & VARCHAR(20) & & & & Versión del pipeline. Defecto: ``1.0.0''. \\
\bottomrule
\caption{Tabla \texttt{dato\_procesado} — Datos procesados (210 registros)}
\end{longtable}

% ══════════════════════════════════════════════════════════════════════════
\section{Tablas del módulo IA — Análisis de Sentimiento}

% ── Tabla: analisis_sentimiento ───────────────────────────────────────────
\subsection{Tabla: analisis\_sentimiento}
\textit{Resultados del modelo de IA para análisis de sentimiento sobre datos procesados. Incluye puntajes de probabilidad para cada clase y nivel de confianza.}

\begin{longtable}{p{4.2cm} p{2.8cm} c c c p{4.5cm}}
\toprule
\tableheadrow
\tablehead{Campo} & \tablehead{Tipo} & \tablehead{PK} & \tablehead{FK} & \tablehead{NN} & \tablehead{Descripción} \\
\midrule
\endhead
id\_analisis & INTEGER & \pk & & & Identificador único del análisis. \\
\rowcolor{tablerow}
id\_dato\_procesado & INTEGER & & \fk & \nn & Ref.\ a \texttt{dato\_procesado}. \\
sentimiento\_predicho & VARCHAR(20) & & & \nn & Resultado: positivo, neutral, negativo. \\
\rowcolor{tablerow}
confianza & DECIMAL(5,4) & & & \nn & Nivel de confianza del modelo (0--1). \\
probabilidad\_positivo & DECIMAL(5,4) & & & & P(positivo). \\
\rowcolor{tablerow}
probabilidad\_neutral & DECIMAL(5,4) & & & & P(neutral). \\
probabilidad\_negativo & DECIMAL(5,4) & & & & P(negativo). \\
\rowcolor{tablerow}
modelo\_version & VARCHAR(50) & & & & Versión del modelo. Defecto: ``1.0.0''. \\
fecha\_analisis & TIMESTAMP & & & & Fecha del análisis. \\
\bottomrule
\caption{Tabla \texttt{analisis\_sentimiento} — Sentimientos (244 registros)}
\end{longtable}

% ── Tabla: analisis_comentario ────────────────────────────────────────────
\subsection{Tabla: analisis\_comentario}
\textit{Sentimiento calculado para comentarios individuales.}

\begin{longtable}{p{4.2cm} p{2.8cm} c c c p{4.5cm}}
\toprule
\tableheadrow
\tablehead{Campo} & \tablehead{Tipo} & \tablehead{PK} & \tablehead{FK} & \tablehead{NN} & \tablehead{Descripción} \\
\midrule
\endhead
id\_analisis & INTEGER & \pk & & & Identificador único. \\
\rowcolor{tablerow}
id\_comentario & INTEGER & & \fk & \nn & Ref.\ a \texttt{comentario}. \\
sentimiento & VARCHAR(20) & & & \nn & Resultado: positivo, neutral, negativo. \\
\rowcolor{tablerow}
confianza & DECIMAL(5,4) & & & \nn & Confianza del modelo. \\
probabilidad\_positivo & DECIMAL(5,4) & & & & P(positivo). \\
\rowcolor{tablerow}
probabilidad\_neutral & DECIMAL(5,4) & & & & P(neutral). \\
probabilidad\_negativo & DECIMAL(5,4) & & & & P(negativo). \\
\rowcolor{tablerow}
fecha\_analisis & TIMESTAMP & & & & Fecha del análisis. \\
\bottomrule
\caption{Tabla \texttt{analisis\_comentario} — Sentimiento de comentarios (448 registros)}
\end{longtable}

% ══════════════════════════════════════════════════════════════════════════
\section{Tablas del módulo NLP / ML}

% ── Tabla: clasificacion_tematica ─────────────────────────────────────────
\subsection{Tabla: clasificacion\_tematica}
\textit{Clasificaciones temáticas: quejas, sugerencias, felicitaciones, información académica, etc.}

\begin{longtable}{p{4.2cm} p{2.8cm} c c c p{4.5cm}}
\toprule
\tableheadrow
\tablehead{Campo} & \tablehead{Tipo} & \tablehead{PK} & \tablehead{FK} & \tablehead{NN} & \tablehead{Descripción} \\
\midrule
\endhead
id & INTEGER & \pk & & & Identificador único. \\
\rowcolor{tablerow}
id\_contenido & INTEGER & & & \nn & ID del contenido analizado. \\
tipo\_contenido & VARCHAR(50) & & & \nn & Tipo: ``post'' o ``comentario''. \\
\rowcolor{tablerow}
tema\_principal & VARCHAR(100) & & & \nn & Tema principal detectado. \\
tema\_secundario & VARCHAR(100) & & & & Tema secundario opcional. \\
\rowcolor{tablerow}
palabras\_clave & TEXT & & & & Keywords asociadas al tema. \\
confianza & DECIMAL(5,4) & & & & Confianza de la clasificación. \\
\rowcolor{tablerow}
es\_academico & BOOLEAN & & & & Si el contenido es académico. \\
es\_relevante\_uebu & BOOLEAN & & & & Si es relevante para la UEBU. \\
\rowcolor{tablerow}
fecha\_clasificacion & TIMESTAMP & & & & Fecha del análisis. \\
\bottomrule
\caption{Tabla \texttt{clasificacion\_tematica} — Clasificación temática (233 registros)}
\end{longtable}

% ── Tabla: nlp_keywords ───────────────────────────────────────────────────
\subsection{Tabla: nlp\_keywords}
\textit{Palabras clave extraídas mediante TF-IDF.}

\begin{longtable}{p{4.2cm} p{2.8cm} c c c p{4.5cm}}
\toprule
\tableheadrow
\tablehead{Campo} & \tablehead{Tipo} & \tablehead{PK} & \tablehead{FK} & \tablehead{NN} & \tablehead{Descripción} \\
\midrule
\endhead
id & INTEGER & \pk & & & Identificador único. \\
\rowcolor{tablerow}
palabra & TEXT & & & \nn & Palabra o término clave. \\
tfidf\_score & REAL & & & & Puntaje TF-IDF. \\
\rowcolor{tablerow}
frecuencia & INTEGER & & & & Frecuencia absoluta en el corpus. \\
tipo & TEXT & & & & Tipo: unigrama, bigrama, entidad. \\
\rowcolor{tablerow}
contexto & TEXT & & & & Fragmentos de texto donde aparece. \\
fecha\_analisis & TIMESTAMP & & & & Fecha del análisis NLP. \\
\bottomrule
\caption{Tabla \texttt{nlp\_keywords} — Palabras clave (50 registros)}
\end{longtable}

% ── Tabla: nlp_topicos ────────────────────────────────────────────────────
\subsection{Tabla: nlp\_topicos}
\textit{Tópicos generados por Latent Dirichlet Allocation (LDA).}

\begin{longtable}{p{4.2cm} p{2.8cm} c c c p{4.5cm}}
\toprule
\tableheadrow
\tablehead{Campo} & \tablehead{Tipo} & \tablehead{PK} & \tablehead{FK} & \tablehead{NN} & \tablehead{Descripción} \\
\midrule
\endhead
id & INTEGER & \pk & & & Identificador único. \\
\rowcolor{tablerow}
metodo & TEXT & & & \nn & Método utilizado (``lda''). \\
topico\_id & INTEGER & & & & Número del tópico. \\
\rowcolor{tablerow}
nombre\_topico & TEXT & & & & Etiqueta asignada al tópico. \\
palabras\_clave & TEXT & & & & Palabras representativas del tópico. \\
\rowcolor{tablerow}
peso\_topico & REAL & & & & Proporción del corpus cubierta. \\
num\_documentos & INTEGER & & & & Documentos asignados al tópico. \\
\rowcolor{tablerow}
coherencia & REAL & & & & Coherencia semántica del tópico. \\
fecha\_analisis & TIMESTAMP & & & & Fecha del análisis. \\
\bottomrule
\caption{Tabla \texttt{nlp\_topicos} — Tópicos LDA (6 registros)}
\end{longtable}

% ── Tabla: nlp_clusters ───────────────────────────────────────────────────
\subsection{Tabla: nlp\_clusters}
\textit{Resultados de agrupamiento K-Means sobre documentos.}

\begin{longtable}{p{4.2cm} p{2.8cm} c c c p{4.5cm}}
\toprule
\tableheadrow
\tablehead{Campo} & \tablehead{Tipo} & \tablehead{PK} & \tablehead{FK} & \tablehead{NN} & \tablehead{Descripción} \\
\midrule
\endhead
id & INTEGER & \pk & & & Identificador único. \\
\rowcolor{tablerow}
cluster\_id & INTEGER & & & & Número del cluster. \\
etiqueta & TEXT & & & & Etiqueta descriptiva del cluster. \\
\rowcolor{tablerow}
palabras\_clave & TEXT & & & & Keywords representativas. \\
num\_documentos & INTEGER & & & & Documentos en el cluster. \\
\rowcolor{tablerow}
sentimiento\_predominante & TEXT & & & & Sentimiento mayoritario. \\
textos\_representativos & TEXT & & & & Ejemplos de textos del cluster. \\
\rowcolor{tablerow}
silhouette\_score & REAL & & & & Métrica de calidad del cluster. \\
fecha\_analisis & TIMESTAMP & & & & Fecha del análisis. \\
\bottomrule
\caption{Tabla \texttt{nlp\_clusters} — Clusters K-Means (8 registros)}
\end{longtable}

% ── Tabla: nlp_entidades ──────────────────────────────────────────────────
\subsection{Tabla: nlp\_entidades}
\textit{Entidades nombradas extraídas por NER (Named Entity Recognition).}

\begin{longtable}{p{4.2cm} p{2.8cm} c c c p{4.5cm}}
\toprule
\tableheadrow
\tablehead{Campo} & \tablehead{Tipo} & \tablehead{PK} & \tablehead{FK} & \tablehead{NN} & \tablehead{Descripción} \\
\midrule
\endhead
id & INTEGER & \pk & & & Identificador único. \\
\rowcolor{tablerow}
texto\_original & TEXT & & & & Texto donde se encontró la entidad. \\
entidad & TEXT & & & & Nombre de la entidad detectada. \\
\rowcolor{tablerow}
tipo\_entidad & TEXT & & & & Tipo: PER, ORG, LOC, MISC. \\
frecuencia & INTEGER & & & & Número de ocurrencias. \\
\rowcolor{tablerow}
contextos & TEXT & & & & Fragmentos contextuales. \\
fecha\_analisis & TIMESTAMP & & & & Fecha del análisis. \\
\bottomrule
\caption{Tabla \texttt{nlp\_entidades} — Entidades NER (28 registros)}
\end{longtable}

% ── Tabla: nlp_resumen_ejecutivo ──────────────────────────────────────────
\subsection{Tabla: nlp\_resumen\_ejecutivo}
\textit{Resúmenes ejecutivos generados automáticamente.}

\begin{longtable}{p{4.2cm} p{2.8cm} c c c p{4.5cm}}
\toprule
\tableheadrow
\tablehead{Campo} & \tablehead{Tipo} & \tablehead{PK} & \tablehead{FK} & \tablehead{NN} & \tablehead{Descripción} \\
\midrule
\endhead
id & INTEGER & \pk & & & Identificador único. \\
\rowcolor{tablerow}
tipo\_resumen & TEXT & & & & Tipo: general, semanal, temático. \\
contenido & TEXT & & & & Texto del resumen. \\
\rowcolor{tablerow}
datos\_soporte & TEXT & & & & Datos JSON de respaldo. \\
fecha\_generacion & TIMESTAMP & & & & Fecha de generación. \\
\bottomrule
\caption{Tabla \texttt{nlp\_resumen\_ejecutivo} — Resúmenes (1 registro)}
\end{longtable}

% ══════════════════════════════════════════════════════════════════════════
\section{Tablas del módulo OSINT — Inteligencia}

\subsection{Tabla: osint\_noticias}
\textit{Noticias sobre la EMI extraídas de medios digitales.}

\begin{longtable}{p{4.2cm} p{2.8cm} c c c p{4.5cm}}
\toprule
\tableheadrow
\tablehead{Campo} & \tablehead{Tipo} & \tablehead{PK} & \tablehead{FK} & \tablehead{NN} & \tablehead{Descripción} \\
\midrule
\endhead
id & INTEGER & \pk & & & Identificador único. \\
\rowcolor{tablerow}
titulo & TEXT & & & \nn & Título de la noticia. \\
resumen & TEXT & & & & Resumen del artículo. \\
\rowcolor{tablerow}
fuente & VARCHAR(200) & & & & Medio de comunicación. \\
url & VARCHAR(500) & & & & Enlace al artículo. \\
\rowcolor{tablerow}
fecha\_publicacion & TIMESTAMP & & & & Fecha de la publicación. \\
fecha\_recoleccion & TIMESTAMP & & & & Fecha de extracción. \\
\rowcolor{tablerow}
termino\_busqueda & VARCHAR(200) & & & & Término usado en la búsqueda. \\
relevancia\_score & DECIMAL(5,4) & & & & Puntaje de relevancia (0--1). \\
\rowcolor{tablerow}
temas\_json & TEXT & & & & Temas en JSON. \\
sentimiento & VARCHAR(20) & & & & Sentimiento detectado. \\
\rowcolor{tablerow}
procesado & BOOLEAN & & & & Si fue procesado. \\
\bottomrule
\caption{Tabla \texttt{osint\_noticias} — Noticias (1 registro)}
\end{longtable}

\subsection{Tabla: osint\_tendencias}
\textit{Datos de Google Trends sobre términos relacionados con la EMI.}

\begin{longtable}{p{4.2cm} p{2.8cm} c c c p{4.5cm}}
\toprule
\tableheadrow
\tablehead{Campo} & \tablehead{Tipo} & \tablehead{PK} & \tablehead{FK} & \tablehead{NN} & \tablehead{Descripción} \\
\midrule
\endhead
id & INTEGER & \pk & & & Identificador único. \\
\rowcolor{tablerow}
termino & VARCHAR(200) & & & \nn & Término de búsqueda. \\
region & VARCHAR(50) & & & & País/región. Defecto: ``Bolivia''. \\
\rowcolor{tablerow}
periodo & VARCHAR(50) & & & & Periodo del dato. \\
valor\_interes & INTEGER & & & & Índice de interés (0--100). \\
\rowcolor{tablerow}
fecha\_dato & DATE & & & & Fecha del dato. \\
fecha\_recoleccion & TIMESTAMP & & & & Fecha de extracción. \\
\rowcolor{tablerow}
tipo & VARCHAR(50) & & & & Tipo: google\_trends. \\
metadata\_json & TEXT & & & & Metadatos JSON. \\
\bottomrule
\caption{Tabla \texttt{osint\_tendencias} — Tendencias (108 registros)}
\end{longtable}

\subsection{Tabla: osint\_fuente\_resumen}
\textit{Resumen de técnicas OSINT aplicadas.}

\begin{longtable}{p{4.2cm} p{2.8cm} c c c p{4.5cm}}
\toprule
\tableheadrow
\tablehead{Campo} & \tablehead{Tipo} & \tablehead{PK} & \tablehead{FK} & \tablehead{NN} & \tablehead{Descripción} \\
\midrule
\endhead
id & INTEGER & \pk & & & Identificador único. \\
\rowcolor{tablerow}
tipo\_tecnica & VARCHAR(50) & & & \nn & Tipo de técnica OSINT. \\
nombre\_fuente & VARCHAR(200) & & & \nn & Nombre de la fuente. \\
\rowcolor{tablerow}
descripcion & TEXT & & & & Descripción de la técnica. \\
url\_base & VARCHAR(500) & & & & URL base de la fuente. \\
\rowcolor{tablerow}
total\_datos\_recolectados & INTEGER & & & & Datos recolectados. \\
ultima\_recoleccion & TIMESTAMP & & & & Última recolección. \\
\rowcolor{tablerow}
estado & VARCHAR(20) & & & & Estado: activo, inactivo. \\
metadata\_json & TEXT & & & & Metadatos JSON. \\
\bottomrule
\caption{Tabla \texttt{osint\_fuente\_resumen} — Resumen de fuentes OSINT (3 registros)}
\end{longtable}

\subsection{Tabla: patron\_identificado}
\textit{Patrones de comportamiento detectados por los algoritmos de IA.}

\begin{longtable}{p{4.2cm} p{2.8cm} c c c p{4.5cm}}
\toprule
\tableheadrow
\tablehead{Campo} & \tablehead{Tipo} & \tablehead{PK} & \tablehead{FK} & \tablehead{NN} & \tablehead{Descripción} \\
\midrule
\endhead
id & INTEGER & \pk & & & Identificador único. \\
\rowcolor{tablerow}
nombre\_patron & VARCHAR(200) & & & \nn & Nombre del patrón. \\
tipo\_patron & VARCHAR(100) & & & \nn & Tipo: temporal, estacional, engagement. \\
\rowcolor{tablerow}
descripcion & TEXT & & & & Descripción del patrón. \\
fuentes\_involucradas & TEXT & & & & Fuentes donde se detectó. \\
\rowcolor{tablerow}
frecuencia & INTEGER & & & & Veces detectado. \\
impacto & VARCHAR(20) & & & & Nivel de impacto. \\
\rowcolor{tablerow}
fecha\_primera\_deteccion & TIMESTAMP & & & & Primera detección. \\
fecha\_ultima\_deteccion & TIMESTAMP & & & & Última detección. \\
\rowcolor{tablerow}
datos\_soporte\_json & TEXT & & & & Datos de soporte. \\
recomendacion\_accion & TEXT & & & & Acción recomendada. \\
\rowcolor{tablerow}
estado & VARCHAR(20) & & & & Estado: activo. \\
relevancia\_vicerrectorado & BOOLEAN & & & & Si es relevante para el Vicerrectorado. \\
\bottomrule
\caption{Tabla \texttt{patron\_identificado} — Patrones (10 registros)}
\end{longtable}

% ══════════════════════════════════════════════════════════════════════════
\section{Tablas del módulo de Usuarios y Seguridad}

% ── Tabla: usuario ────────────────────────────────────────────────────────
\subsection{Tabla: usuario}
\textit{Gestión de roles, emails y estado de los usuarios del sistema.}

\begin{longtable}{p{4.2cm} p{2.8cm} c c c p{4.5cm}}
\toprule
\tableheadrow
\tablehead{Campo} & \tablehead{Tipo} & \tablehead{PK} & \tablehead{FK} & \tablehead{NN} & \tablehead{Descripción} \\
\midrule
\endhead
id\_usuario & INTEGER & \pk & & & Identificador único. \\
\rowcolor{tablerow}
username & TEXT & & & \nn & Nombre de usuario único. \\
email & TEXT & & & \nn & Correo electrónico institucional. \\
\rowcolor{tablerow}
password\_hash & TEXT & & & \nn & Hash SHA-256 de la contraseña. \\
nombre\_completo & TEXT & & & \nn & Nombre completo del usuario. \\
\rowcolor{tablerow}
rol & TEXT & & & \nn & Rol: administrador, vicerrector, uebu. \\
cargo & TEXT & & & & Cargo institucional. \\
\rowcolor{tablerow}
activo & INTEGER & & & & Estado activo (1) o inactivo (0). \\
ultimo\_login & TEXT & & & & Fecha del último inicio de sesión. \\
\rowcolor{tablerow}
fecha\_creacion & TEXT & & & & Fecha de creación de la cuenta. \\
permisos & TEXT & & & & Permisos JSON del usuario. Defecto: ``\{\}''. \\
\bottomrule
\caption{Tabla \texttt{usuario} — Usuarios del sistema (6 registros)}
\end{longtable}

\subsection{Tabla: log\_actividad}
\textit{Registro de acciones realizadas por los usuarios en el sistema.}

\begin{longtable}{p{4.2cm} p{2.8cm} c c c p{4.5cm}}
\toprule
\tableheadrow
\tablehead{Campo} & \tablehead{Tipo} & \tablehead{PK} & \tablehead{FK} & \tablehead{NN} & \tablehead{Descripción} \\
\midrule
\endhead
id\_log & INTEGER & \pk & & & Identificador único. \\
\rowcolor{tablerow}
id\_usuario & INTEGER & & \fk & & Ref.\ a \texttt{usuario}. \\
accion & TEXT & & & \nn & Tipo de acción realizada. \\
\rowcolor{tablerow}
detalle & TEXT & & & & Descripción detallada. \\
ip\_address & TEXT & & & & Dirección IP del usuario. \\
\rowcolor{tablerow}
fecha & TEXT & & & & Fecha y hora de la acción. \\
\bottomrule
\caption{Tabla \texttt{log\_actividad} — Logs de actividad (47 registros)}
\end{longtable}

% ══════════════════════════════════════════════════════════════════════════
\section{Tablas del módulo de Alertas}

\subsection{Tabla: alerta}
\textit{Alertas generadas automáticamente por el sistema ante situaciones críticas.}

\begin{longtable}{p{4.2cm} p{2.8cm} c c c p{4.5cm}}
\toprule
\tableheadrow
\tablehead{Campo} & \tablehead{Tipo} & \tablehead{PK} & \tablehead{FK} & \tablehead{NN} & \tablehead{Descripción} \\
\midrule
\endhead
id\_alerta & INTEGER & \pk & & & Identificador único. \\
\rowcolor{tablerow}
tipo & TEXT & & & \nn & Tipo de alerta. \\
severidad & TEXT & & & \nn & Niveles: critica, alta, media, baja. \\
\rowcolor{tablerow}
titulo & TEXT & & & \nn & Título descriptivo de la alerta. \\
descripcion & TEXT & & & & Descripción detallada. \\
\rowcolor{tablerow}
fuente & TEXT & & & & Fuente que originó la alerta. Defecto: ``facebook''. \\
estado & TEXT & & & \nn & Estado: nueva, en\_proceso, resuelta. \\
\rowcolor{tablerow}
id\_dato\_procesado & INTEGER & & & & Referencia al dato que activó la alerta. \\
confianza & REAL & & & & Confianza de la detección. \\
\rowcolor{tablerow}
engagement & INTEGER & & & & Engagement asociado. \\
asignado\_a & INTEGER & & \fk & & Ref.\ a \texttt{usuario} asignado. \\
\rowcolor{tablerow}
resolucion & TEXT & & & & Notas de resolución. \\
resuelto\_por & INTEGER & & \fk & & Ref.\ a \texttt{usuario} que resolvió. \\
\rowcolor{tablerow}
fecha\_creacion & TEXT & & & & Fecha de creación. \\
fecha\_resolucion & TEXT & & & & Fecha de resolución. \\
\bottomrule
\caption{Tabla \texttt{alerta} — Alertas (20 registros)}
\end{longtable}

\subsection{Tabla: configuracion\_alerta}
\textit{Reglas configurables para la generación automática de alertas.}

\begin{longtable}{p{4.2cm} p{2.8cm} c c c p{4.5cm}}
\toprule
\tableheadrow
\tablehead{Campo} & \tablehead{Tipo} & \tablehead{PK} & \tablehead{FK} & \tablehead{NN} & \tablehead{Descripción} \\
\midrule
\endhead
id\_config & INTEGER & \pk & & & Identificador único. \\
\rowcolor{tablerow}
nombre & TEXT & & & \nn & Nombre de la regla. \\
tipo\_alerta & TEXT & & & \nn & Tipo de alerta que genera. \\
\rowcolor{tablerow}
umbral\_valor & REAL & & & & Valor umbral de activación. Defecto: 0.7. \\
umbral\_confianza & REAL & & & & Confianza mínima. Defecto: 0.5. \\
\rowcolor{tablerow}
severidad\_minima & TEXT & & & & Severidad mínima. Defecto: ``media''. \\
activa & INTEGER & & & & Si la regla está activa. Defecto: 1. \\
\rowcolor{tablerow}
notificar\_email & INTEGER & & & & Si envía email. Defecto: 0. \\
creado\_por & INTEGER & & \fk & & Ref.\ a \texttt{usuario} que creó la regla. \\
\rowcolor{tablerow}
fecha\_creacion & TEXT & & & & Fecha de creación. \\
\bottomrule
\caption{Tabla \texttt{configuracion\_alerta} — Configuraciones de alerta (4 registros)}
\end{longtable}

% ══════════════════════════════════════════════════════════════════════════
\section{Tablas de Evaluación y Auditoría}

\subsection{Tabla: evaluacion\_sistema}
\textit{Métricas de evaluación de calidad y rendimiento del sistema.}

\begin{longtable}{p{4.2cm} p{2.8cm} c c c p{4.5cm}}
\toprule
\tableheadrow
\tablehead{Campo} & \tablehead{Tipo} & \tablehead{PK} & \tablehead{FK} & \tablehead{NN} & \tablehead{Descripción} \\
\midrule
\endhead
id & INTEGER & \pk & & & Identificador único. \\
\rowcolor{tablerow}
metrica & TEXT & & & \nn & Nombre de la métrica. \\
categoria & TEXT & & & & Categoría: recolección, sentimiento, nlp\_ml, rendimiento. \\
\rowcolor{tablerow}
valor & REAL & & & & Valor de la métrica (0--100). \\
detalle & TEXT & & & & Descripción detallada. \\
\rowcolor{tablerow}
fecha\_evaluacion & TIMESTAMP & & & & Fecha de evaluación. \\
\bottomrule
\caption{Tabla \texttt{evaluacion\_sistema} — Evaluación del sistema (22 registros)}
\end{longtable}

\subsection{Tabla: log\_ejecucion}
\textit{Historial de ejecuciones de los procesos de recolección y análisis.}

\begin{longtable}{p{4.2cm} p{2.8cm} c c c p{4.5cm}}
\toprule
\tableheadrow
\tablehead{Campo} & \tablehead{Tipo} & \tablehead{PK} & \tablehead{FK} & \tablehead{NN} & \tablehead{Descripción} \\
\midrule
\endhead
id\_log & INTEGER & \pk & & & Identificador único. \\
\rowcolor{tablerow}
tipo\_operacion & VARCHAR(50) & & & \nn & Tipo: scraping, etl, nlp, evaluación. \\
fuente & VARCHAR(100) & & & & Fuente asociada. \\
\rowcolor{tablerow}
fecha\_inicio & TIMESTAMP & & & \nn & Inicio de la ejecución. \\
fecha\_fin & TIMESTAMP & & & & Fin de la ejecución. \\
\rowcolor{tablerow}
duracion\_segundos & DECIMAL(10,2) & & & & Duración en segundos. \\
registros\_procesados & INTEGER & & & & Total de registros procesados. \\
\rowcolor{tablerow}
registros\_exitosos & INTEGER & & & & Registros procesados con éxito. \\
registros\_fallidos & INTEGER & & & & Registros fallidos. \\
\rowcolor{tablerow}
estado & VARCHAR(20) & & & & Estado: completado, error. \\
mensaje\_error & TEXT & & & & Mensaje de error si aplica. \\
\rowcolor{tablerow}
detalles\_json & TEXT & & & & Detalles adicionales JSON. \\
\bottomrule
\caption{Tabla \texttt{log\_ejecucion} — Logs de ejecución (46 registros)}
\end{longtable}

% ══════════════════════════════════════════════════════════════════════════
\section{Tablas del módulo de IA Avanzada}

\subsection{Tabla: alerta\_ia}
\textit{Alertas generadas por los módulos de Inteligencia Artificial del sistema.}

\begin{longtable}{p{4.2cm} p{2.8cm} c c c p{4.5cm}}
\toprule
\tableheadrow
\tablehead{Campo} & \tablehead{Tipo} & \tablehead{PK} & \tablehead{FK} & \tablehead{NN} & \tablehead{Descripción} \\
\midrule
\endhead
id\_alerta & INTEGER & \pk & & & Identificador único. \\
\rowcolor{tablerow}
tipo\_alerta & VARCHAR(50) & & & \nn & Tipo de alerta IA. \\
titulo & VARCHAR(200) & & & \nn & Título descriptivo. \\
\rowcolor{tablerow}
descripcion & TEXT & & & \nn & Descripción detallada. \\
severidad & VARCHAR(20) & & & \nn & Niveles: critica, alta, media, baja. \\
\rowcolor{tablerow}
modulo\_origen & VARCHAR(50) & & & & Módulo que generó la alerta. \\
referencia\_id & INTEGER & & & & ID del registro relacionado. \\
\rowcolor{tablerow}
referencia\_tipo & VARCHAR(50) & & & & Tipo de registro referenciado. \\
requiere\_accion & BOOLEAN & & & & Si requiere acción inmediata. Defecto: 0. \\
\rowcolor{tablerow}
estado & VARCHAR(20) & & & & Estado: nueva, leida, resuelta. Defecto: ``nueva''. \\
fecha\_creacion & TIMESTAMP & & & & Fecha de creación. \\
\rowcolor{tablerow}
fecha\_lectura & TIMESTAMP & & & & Fecha de lectura. \\
fecha\_resolucion & TIMESTAMP & & & & Fecha de resolución. \\
\rowcolor{tablerow}
usuario\_asignado & VARCHAR(100) & & & & Usuario asignado a resolver. \\
notas\_resolucion & TEXT & & & & Notas de la resolución. \\
\bottomrule
\caption{Tabla \texttt{alerta\_ia} — Alertas de IA (0 registros)}
\end{longtable}

\subsection{Tabla: anomalia\_detectada}
\textit{Anomalías detectadas por el algoritmo Isolation Forest sobre métricas del sistema.}

\begin{longtable}{p{4.2cm} p{2.8cm} c c c p{4.5cm}}
\toprule
\tableheadrow
\tablehead{Campo} & \tablehead{Tipo} & \tablehead{PK} & \tablehead{FK} & \tablehead{NN} & \tablehead{Descripción} \\
\midrule
\endhead
id\_anomalia & INTEGER & \pk & & & Identificador único. \\
\rowcolor{tablerow}
tipo\_anomalia & VARCHAR(50) & & & \nn & Tipo de anomalía detectada. \\
descripcion & TEXT & & & \nn & Descripción de la anomalía. \\
\rowcolor{tablerow}
severidad & VARCHAR(20) & & & \nn & Niveles: critica, alta, media, baja. \\
metrica\_afectada & VARCHAR(100) & & & & Métrica afectada. \\
\rowcolor{tablerow}
valor\_esperado & DECIMAL(15,4) & & & & Valor esperado por el modelo. \\
valor\_observado & DECIMAL(15,4) & & & & Valor observado (anómalo). \\
\rowcolor{tablerow}
anomaly\_score & DECIMAL(10,6) & & & & Puntuación de anomalía (Isolation Forest). \\
fecha\_deteccion & TIMESTAMP & & & & Fecha de detección. \\
\rowcolor{tablerow}
fecha\_ocurrencia & TIMESTAMP & & & & Fecha real de la ocurrencia. \\
estado & VARCHAR(20) & & & & Estado: nueva, investigada, resuelta. Defecto: ``nueva''. \\
\rowcolor{tablerow}
notas & TEXT & & & & Notas adicionales. \\
metadata\_json & TEXT & & & & Metadatos JSON. \\
\bottomrule
\caption{Tabla \texttt{anomalia\_detectada} — Anomalías (5 registros)}
\end{longtable}

\subsection{Tabla: modelo\_entrenamiento}
\textit{Registro de modelos de Machine Learning entrenados y sus métricas de rendimiento.}

\begin{longtable}{p{4.2cm} p{2.8cm} c c c p{4.5cm}}
\toprule
\tableheadrow
\tablehead{Campo} & \tablehead{Tipo} & \tablehead{PK} & \tablehead{FK} & \tablehead{NN} & \tablehead{Descripción} \\
\midrule
\endhead
id\_entrenamiento & INTEGER & \pk & & & Identificador único. \\
\rowcolor{tablerow}
tipo\_modelo & VARCHAR(50) & & & \nn & Tipo: sentiment, clustering, topic. \\
nombre\_modelo & VARCHAR(100) & & & \nn & Nombre del modelo. \\
\rowcolor{tablerow}
version & VARCHAR(20) & & & \nn & Versión del modelo. \\
metricas\_json & TEXT & & & \nn & Métricas de rendimiento (accuracy, f1, etc.). \\
\rowcolor{tablerow}
parametros\_json & TEXT & & & & Hiperparámetros del modelo. \\
datos\_entrenamiento & INTEGER & & & & Registros usados para entrenar. \\
\rowcolor{tablerow}
datos\_validacion & INTEGER & & & & Registros usados para validar. \\
modelo\_path & VARCHAR(500) & & & & Ruta al archivo del modelo. \\
\rowcolor{tablerow}
estado & VARCHAR(20) & & & & Estado: entrenado, activo, deprecado. \\
fecha\_entrenamiento & TIMESTAMP & & & & Fecha de entrenamiento. \\
\rowcolor{tablerow}
fecha\_activacion & TIMESTAMP & & & & Fecha de activación en producción. \\
\bottomrule
\caption{Tabla \texttt{modelo\_entrenamiento} — Modelos ML (1 registro)}
\end{longtable}

\subsection{Tabla: clustering\_modelo}
\textit{Modelos de clustering K-Means con métricas de calidad (silhouette, Calinski-Harabasz, Davies-Bouldin).}

\begin{longtable}{p{4.2cm} p{2.8cm} c c c p{4.5cm}}
\toprule
\tableheadrow
\tablehead{Campo} & \tablehead{Tipo} & \tablehead{PK} & \tablehead{FK} & \tablehead{NN} & \tablehead{Descripción} \\
\midrule
\endhead
id\_modelo & INTEGER & \pk & & & Identificador único. \\
\rowcolor{tablerow}
nombre\_modelo & VARCHAR(100) & & & \nn & Nombre del modelo de clustering. \\
n\_clusters & INTEGER & & & \nn & Número de clusters (k). \\
\rowcolor{tablerow}
silhouette\_score & DECIMAL(5,4) & & & & Coeficiente de Silhouette (0--1). \\
calinski\_score & DECIMAL(10,2) & & & & Índice Calinski-Harabasz. \\
\rowcolor{tablerow}
davies\_bouldin\_score & DECIMAL(10,4) & & & & Índice Davies-Bouldin. \\
total\_documentos & INTEGER & & & & Documentos procesados. \\
\rowcolor{tablerow}
cluster\_keywords\_json & TEXT & & & & Keywords por cluster (JSON). \\
configuracion\_json & TEXT & & & & Configuración del modelo (JSON). \\
\rowcolor{tablerow}
modelo\_path & VARCHAR(500) & & & & Ruta al archivo del modelo. \\
activo & BOOLEAN & & & & Si el modelo está activo. Defecto: 1. \\
\rowcolor{tablerow}
fecha\_entrenamiento & TIMESTAMP & & & & Fecha de entrenamiento. \\
\bottomrule
\caption{Tabla \texttt{clustering\_modelo} — Modelos de clustering (0 registros)}
\end{longtable}

\subsection{Tabla: clustering\_resultado}
\textit{Asignación de cada documento procesado a un cluster específico.}

\begin{longtable}{p{4.2cm} p{2.8cm} c c c p{4.5cm}}
\toprule
\tableheadrow
\tablehead{Campo} & \tablehead{Tipo} & \tablehead{PK} & \tablehead{FK} & \tablehead{NN} & \tablehead{Descripción} \\
\midrule
\endhead
id\_clustering & INTEGER & \pk & & & Identificador único. \\
\rowcolor{tablerow}
id\_dato\_procesado & INTEGER & & \fk & \nn & Ref.\ a \texttt{dato\_procesado}. \\
cluster\_id & INTEGER & & & \nn & ID del cluster asignado. \\
\rowcolor{tablerow}
distancia\_centroide & DECIMAL(10,6) & & & & Distancia al centroide del cluster. \\
confianza\_asignacion & DECIMAL(5,4) & & & & Confianza de la asignación. \\
\rowcolor{tablerow}
fecha\_clustering & TIMESTAMP & & & & Fecha del clustering. \\
\bottomrule
\caption{Tabla \texttt{clustering\_resultado} — Resultados de clustering (0 registros)}
\end{longtable}

\subsection{Tabla: analisis\_tendencia}
\textit{Análisis de tendencias temporales sobre las métricas del sistema.}

\begin{longtable}{p{4.2cm} p{2.8cm} c c c p{4.5cm}}
\toprule
\tableheadrow
\tablehead{Campo} & \tablehead{Tipo} & \tablehead{PK} & \tablehead{FK} & \tablehead{NN} & \tablehead{Descripción} \\
\midrule
\endhead
id\_tendencia & INTEGER & \pk & & & Identificador único. \\
\rowcolor{tablerow}
periodo & VARCHAR(50) & & & \nn & Periodo analizado. \\
metrica & VARCHAR(100) & & & \nn & Métrica evaluada. \\
\rowcolor{tablerow}
tipo\_tendencia & VARCHAR(30) & & & \nn & Tipo: creciente, decreciente, estable. \\
valor\_slope & DECIMAL(10,6) & & & & Pendiente de la regresión. \\
\rowcolor{tablerow}
valor\_r\_squared & DECIMAL(5,4) & & & & Coeficiente de determinación $R^2$. \\
confianza & DECIMAL(5,4) & & & & Confianza del análisis. \\
\rowcolor{tablerow}
fecha\_inicio & DATE & & & \nn & Inicio del periodo. \\
fecha\_fin & DATE & & & \nn & Fin del periodo. \\
\rowcolor{tablerow}
datos\_puntos & INTEGER & & & & Puntos de datos utilizados. \\
estacionalidad\_detectada & BOOLEAN & & & & Si se detectó estacionalidad. \\
\rowcolor{tablerow}
forecast\_json & TEXT & & & & Pronóstico generado (JSON). \\
metadata\_json & TEXT & & & & Metadatos adicionales. \\
\rowcolor{tablerow}
fecha\_analisis & TIMESTAMP & & & & Fecha del análisis. \\
\bottomrule
\caption{Tabla \texttt{analisis\_tendencia} — Tendencias (5 registros)}
\end{longtable}

\subsection{Tabla: anotacion\_manual}
\textit{Anotaciones manuales de sentimiento para validación (gold standard) de los modelos de IA.}

\begin{longtable}{p{4.2cm} p{2.8cm} c c c p{4.5cm}}
\toprule
\tableheadrow
\tablehead{Campo} & \tablehead{Tipo} & \tablehead{PK} & \tablehead{FK} & \tablehead{NN} & \tablehead{Descripción} \\
\midrule
\endhead
id\_anotacion & INTEGER & \pk & & & Identificador único. \\
\rowcolor{tablerow}
id\_dato\_procesado & INTEGER & & \fk & \nn & Ref.\ a \texttt{dato\_procesado}. \\
texto\_original & TEXT & & & \nn & Texto que fue anotado. \\
\rowcolor{tablerow}
sentimiento\_anotado & VARCHAR(20) & & & \nn & Sentimiento asignado manualmente. \\
confianza\_anotacion & VARCHAR(20) & & & & Nivel de confianza del anotador. \\
\rowcolor{tablerow}
anotador & VARCHAR(100) & & & & Nombre del anotador. \\
fecha\_anotacion & TIMESTAMP & & & & Fecha de la anotación. \\
\rowcolor{tablerow}
notas & TEXT & & & & Notas del anotador. \\
usado\_entrenamiento & BOOLEAN & & & & Si fue usado para entrenar. Defecto: 0. \\
\bottomrule
\caption{Tabla \texttt{anotacion\_manual} — Anotaciones manuales (1 registro)}
\end{longtable}

\subsection{Tabla: correlacion\_resultado}
\textit{Resultados de análisis de correlación estadística entre variables del sistema.}

\begin{longtable}{p{4.2cm} p{2.8cm} c c c p{4.5cm}}
\toprule
\tableheadrow
\tablehead{Campo} & \tablehead{Tipo} & \tablehead{PK} & \tablehead{FK} & \tablehead{NN} & \tablehead{Descripción} \\
\midrule
\endhead
id\_correlacion & INTEGER & \pk & & & Identificador único. \\
\rowcolor{tablerow}
variable\_1 & VARCHAR(100) & & & \nn & Primera variable. \\
variable\_2 & VARCHAR(100) & & & \nn & Segunda variable. \\
\rowcolor{tablerow}
coeficiente\_correlacion & DECIMAL(6,4) & & & \nn & Coeficiente de correlación. \\
p\_value & DECIMAL(10,8) & & & \nn & Valor p (significancia). \\
\rowcolor{tablerow}
es\_significativa & BOOLEAN & & & \nn & Si la correlación es significativa ($p < 0.05$). \\
fuerza & VARCHAR(30) & & & & Fuerza: débil, moderada, fuerte. \\
\rowcolor{tablerow}
direccion & VARCHAR(20) & & & & Dirección: positiva, negativa. \\
n\_muestras & INTEGER & & & & Número de muestras. \\
\rowcolor{tablerow}
metodo & VARCHAR(30) & & & & Método: pearson, spearman. Defecto: ``pearson''. \\
fecha\_analisis & TIMESTAMP & & & & Fecha del análisis. \\
\bottomrule
\caption{Tabla \texttt{correlacion\_resultado} — Correlaciones (28 registros)}
\end{longtable}

% ══════════════════════════════════════════════════════════════════════════
\section{Resumen estadístico}

\begin{table}[H]
\centering
\begin{tabularx}{\textwidth}{l r X}
\toprule
\textbf{Tabla} & \textbf{Registros} & \textbf{Módulo} \\
\midrule
fuente\_osint & 4 & Recolección OSINT \\
dato\_recolectado & 214 & Recolección OSINT \\
comentario & 534 & Recolección OSINT \\
dato\_procesado & 210 & ETL / Procesamiento \\
analisis\_sentimiento & 244 & IA — Sentimiento \\
analisis\_comentario & 448 & IA — Sentimiento \\
clasificacion\_tematica & 233 & NLP / ML \\
nlp\_keywords & 50 & NLP / ML \\
nlp\_topicos & 6 & NLP / ML \\
nlp\_clusters & 8 & NLP / ML \\
nlp\_entidades & 28 & NLP / ML \\
nlp\_resumen\_ejecutivo & 1 & NLP / ML \\
osint\_noticias & 1 & Inteligencia OSINT \\
osint\_tendencias & 108 & Inteligencia OSINT \\
osint\_fuente\_resumen & 3 & Inteligencia OSINT \\
patron\_identificado & 10 & Inteligencia OSINT \\
usuario & 6 & Seguridad \\
log\_actividad & 47 & Seguridad \\
alerta & 20 & Alertas \\
alerta\_ia & 0 & Alertas IA \\
configuracion\_alerta & 4 & Alertas \\
anomalia\_detectada & 5 & IA Avanzada \\
modelo\_entrenamiento & 1 & IA Avanzada \\
clustering\_modelo & 0 & IA Avanzada \\
clustering\_resultado & 0 & IA Avanzada \\
analisis\_tendencia & 5 & IA Avanzada \\
anotacion\_manual & 1 & IA Avanzada \\
correlacion\_resultado & 28 & IA Avanzada \\
evaluacion\_sistema & 22 & Evaluación \\
log\_ejecucion & 46 & Auditoría \\
\midrule
\textbf{Total} & \textbf{2287} & \textbf{30 tablas} \\
\bottomrule
\end{tabularx}
\caption{Resumen de registros por tabla}
\end{table}

% ══════════════════════════════════════════════════════════════════════════
\newpage
\section*{Control de versiones del documento}
\begin{table}[H]
\centering
\begin{tabularx}{\textwidth}{c c l X}
\toprule
\textbf{Versión} & \textbf{Fecha} & \textbf{Autor} & \textbf{Cambios} \\
\midrule
1.0 & Feb 2026 & Alvaro Santiago Encinas Flores & Versión inicial del diccionario de datos. \\
2.0 & May 2026 & Alvaro Santiago Encinas Flores & Sincronización con BD real: 30 tablas, conteos actualizados, 8 tablas agregadas (IA Avanzada). \\
\bottomrule
\end{tabularx}
\end{table}

\end{document}
